% Faithful transcription — original French. Transcribed from the original first printing:
% H. G. Zeuthen, "Études géométriques de quelques-unes des propriétés de deux surfaces
% dont les points se correspondent un-à-un", Mathematische Annalen, Band 4, Heft 1 (1871),
% pp. 21–49. Source: Internet Archive, sim_mathematische-annalen_1871_4 (printed pp. 21–49 =
% leaves n24–n52).
%
% \origpage uses the PRINTED page numbers (21 to 49). Running headers and sheet signature marks
% are omitted per house style. Footnotes are kept at the foot of their respective pages.
% See notation.md for full notation decisions and conventions.
\documentclass{article}
\usepackage{readmasters}
\begin{document}

\origpage{21}
\section*{Études géométriques de quelques-unes des propriétés de deux surfaces dont les points se correspondent un-à-un.}

\subsection*{Par H.~G. \textsc{Zeuthen} à Copenhague.}

\section*{I. Remarques préliminaires${}^{*)}$.}

\textbf{1.} Nous appellerons les deux surfaces $[F_1]$ et $[F_2]$, et nous donnerons aux nombres de leurs singularités les notations indiquées au commencement du précédent article, en y ajoutant seulement, pour distinguer celles des deux surfaces, les suffixes 1 et 2. Nous désignerons encore par $p_1$ et $p_2$ les genres des sections planes de $[F_1]$ et de $[F_2]$, de façon que pour chacune des surfaces
\[
2(p - 1) = n(n - 3) - 2b - 2c = a + c - 2n. \tag{1}
\]

\textbf{2.} Aux deux points de l'une des surfaces qui coïncident dans un même point de sa courbe double, correspondent en général, sur l'autre surface, deux points différents, et le lieu de ces points est une courbe qui correspond à la courbe double. Nous désignerons la courbe double de $[F_1]$ par $(B_1)$, la courbe qui y correspond sur $[F_2]$ par $(B_{1, 2})$ et l'ordre de cette dernière courbe par $b_{1, 2}$, et nous appliquerons les notations analogues $(B_{2, 1})$ et $b_{2, 1}$ à la courbe de $[F_1]$ qui correspond à la courbe double $(B_2)$ de $[F_2]$. --- S'il y a sur les deux surfaces des courbes doubles qui se correspondent, deux branches de la courbe $(B_{1, 2})$ coïncideront dans l'une et deux branches de $(B_{2, 1})$ coïncideront dans l'autre de ces deux courbes doubles correspondantes.

\textbf{3.} \emph{Courbes cuspidales} (3---5). A une courbe cuspidale de l'une des surfaces peut correspondre sur l'autre ou une courbe simple ou une courbe cuspidale, mais nous supposerons que seulement le premier de ces deux cas a lieu pour nos surfaces $[F_1]$ et $[F_2]$, et nous désignerons

\textbf{*)} Nous renvoyons pour une partie de ces remarques préliminaires aux démonstrations analytiques dans un mémoire de Mr.~Nöther, Mathematische Annalen, 2.~Bd.\ 294---298.

\origpage{22}
par $(C_{1, 2})$ la courbe simple de $[F_2]$ qui correspond à la courbe cuspidale $(C_1)$ de $[F_1]$, et par $(C_{2, 1})$ la courbe simple de $[F_1]$ qui correspond à la courbe cuspidale $(C_2)$ de $[F_2]$. Les ordres de ces deux courbes seront nommés $c_{1, 2}$ et $c_{2, 1}$.

Pour nous rendre compte des propriétés de la surface $[F_2]$ à un point de $(C_{1, 2})$, nous regarderons la courbe cuspidale $(C_1)$ à laquelle elle correspond comme limite d'une courbe double $(B_1)$. Si nous désignons par $P_{1, 2}$ et $P'_{1, 2}$ les deux points de la courbe correspondante $(B_{1, 2})$ qui correspondent aux deux points de $[F_1]$ qui se trouvent dans un même point $P_1$ de la courbe double $(B_1)$, on voit qu'une courbe sur $[F_2]$ qui passe par tous les deux points $P_{1, 2}$ et $P'_{1, 2}$ correspond à une courbe sur $[F_1]$ qui a $P_1$ pour point double, mais que les courbes sur $[F_2]$ qui passent seulement par un des points $P_{1, 2}$ et $P'_{1, 2}$ correspondent à des courbes dont une seule branche passe par $P_1$. Si la courbe double $(B_1)$ tend à devenir cuspidale, les points $P_{1, 2}$ et $P'_{1, 2}$ tendront à coïncider. La courbe $(C_{1, 2})$ qui correspond à la courbe cuspidale $(C_1)$ sera donc la limite de deux branches coïncidentes de $(B_{1, 2})$, et comme la droite $(P_{1, 2} P'_{1, 2})$ qui joignait les deux points de $(B_{1, 2})$ qui correspondaient à un même point de $(B_1)$ doit avoir une limite déterminée, il y aura par chaque point $P_{1, 2}$ de $(C_{1, 2})$ une direction qui joint les deux points infiniment voisins qui correspondent à un même point de la courbe cuspidale $(C_1)$. Les autres directions par $P_{1, 2}$ joignent seulement des points des deux branches de $(B_{1, 2})$, coïncidant dans $(C_{1, 2})$, qui correspondent à des points de $(C_1)$ infiniment voisins l'un de l'autre. A une courbe sur $[F_2]$ qui à $P_{1, 2}$ est tangente à la première direction $(P_{1, 2} P'_{1, 2})$ correspond sur $[F_1]$ une courbe douée \emph{d'un point cuspidal au point $P_1$}; mais aux autres courbes qui sur $[F_2]$ passent par $P_{1, 2}$ correspondent des courbes tangentes à $(C_1)$.

\textbf{4.} En regardant une courbe sur $[F_2]$ tangente à $(P_{1, 2} P'_{1, 2})$ et la courbe correspondante sur $[F_1]$, on voit que les parties de $[F_2]$ qui se trouvent des deux côtés de $(C_{1, 2})$ correspondent respectivement aux deux nappes de $[F_1]$ qui sont réunies à la courbe cuspidale $(C_1)$, et que, par conséquent, aussi les courbes par $P_1$ et sur $[F_1]$ dont les correspondantes ne sont pas tangentes à $(P_{1, 2} P'_{1, 2})$, passent en général de l'une des deux nappes à l'autre${}^{*)}$. Une courbe par $P_1$ et sur $[F_1]$ ne reste sur la même nappe que dans le cas où la correspondante sur

\textbf{*)} En général, une courbe qui se trouve sur une surface douée d'une courbe cuspidale $(C)$, passe de l'une des nappes à l'autre au cas où elle a un contact biponctuel avec $(C)$, et reste sur la même nappe au cas d'un contact quatreponctuel, ce qu'on peut prouver ou directement ou au moyen de la représentation dont nous nous occupons.

\origpage{23}
$[F_2]$ est tangente à $(C_{1, 2})$ --- ou a $P_{1, 2}$ pour point cuspidal, --- et alors elle aura avec $(C_1)$ un contact \emph{quatreponctuel}.

Ces considérations nous montrent que les courbes sur $[F_1]$ qui passent par $P_1$, ont, au moins, un contact triponctuel (osculation) avec \emph{une surface qui est tangente à $[F_1]$ le long de la courbe cuspidale $(C_1)$} ou qui seulement au point $P_1$ est osculatrice à $(C_1)$ et tangente à $[F_1]$. Désignons par $[G]$ une surface satisfaisant à ces conditions.

\textbf{5.} A l'exception des courbes tangentes à la direction particulière $(P_{1, 2} P'_{1, 2})$, toute courbe par le point $P_{1, 2}$ et sur $[F_2]$ a donc pour correspondante sur $[F_1]$ une courbe tangente à $(C_1)$ à un même point $P_1$ et osculatrice à la surface $[G]$. Les deux courbes sur $[F_1]$ qui correspondent à deux courbes par $P_{1, 2}$ auront donc, en général, deux points communs coïncidant au point $P_1$. Elles en auront un troisième au cas où les deux courbes sur $[F_2]$ se touchent à $P_{1, 2}$. Par conséquent, à deux courbes sur $[F_2]$ qui ont un contact au point $P_{1, 2}$, correspondent sur $[F_1]$ deux courbes qui ont au point $P_1$ le même plan osculateur. --- La proposition réciproque ne serait pas vraie; car bien que deux courbes sur $[F_1]$ qui ont au point $P_1$ le même plan osculateur aient aussi la même courbure (celle de la section faite par le plan osculateur à la surface $[G]$), et qu'elles aient par conséquent trois points communs coïncidents au lieu de deux, le troisième peut sur les deux courbes appartenir à des nappes différentes. On voit donc que les courbes par $P_1$ et sur $[F_1]$ qui ont le même plan osculateur, n'ont pour correspondantes des courbes tangentes à une même direction par $P_{1, 2}$ que dans le cas où leurs arcs situés du même côté de $P_1$ se trouvent sur la même nappe. Par conséquent, les courbes sur $[F_2]$ dont les correspondantes ont le même plan osculateur au point $P_1$, sont tangentes à l'une ou à l'autre de \emph{deux directions par $P_{1, 2}$}. Le faisceau formé par ces couples de directions est en \emph{involution}.

Les directions d'une même couple ne coïncident que dans les cas où le plan osculateur, qui passe toujours par la tangente de la courbe cuspidale $(C_1)$ au point $P_1$, y est \emph{ou} osculateur à la courbe $(C_1)$ \emph{ou} tangent à la surface $[F_1]$. Elles coïncident au premier cas avec la tangente à la courbe $(C_{1, 2})$ et, au second, avec la direction particulière $(P_{1, 2} P'_{1, 2})$${}^{*)}$.

En appliquant ces considérations à \emph{une section plane faite à la surface $[F_1]$ par un plan passant par la tangente de la courbe $(C_1)$ au}

\textbf{*)} Ces propriétés subiront des altérations lorsque le plan tangent de $[F_1]$ au point $P_1$ est aussi osculateur à la courbe $(C_1)$. Alors les différentes directions par $P_{1, 2}$ correspondront aux différentes courbures des courbes sur $[F_1]$, qui auront le même plan osculateur. --- Un autre cas exceptionnel se présentera lorsque le point $P_{1, 2}$ est un des points, dits fondamentaux, dont nous allons parler dans le \no 6. Alors toutes les courbes par $P_1$ et sur $[F_1]$ auront pour correspondantes des courbes ayant une tangente commune au point $P_{1, 2}$.

\origpage{24}
point $P_1$, on voit que la courbe correspondante sur $[F_2]$ a, en général, \emph{un point double au point $P_{1, 2}$}, que ce point double devient cuspidal au cas où le plan est osculateur à $(C_1)$ ou tangent à $[F_1]$, et que dans ces cas la tangente au point cuspidal est tangente, respectivement, à la courbe $(C_{1, 2})$ ou à la direction $(P_{1, 2} P'_{1, 2})$.

Il va sans dire que la surface $[F_1]$ a, à un point de $(C_{2, 1})$, des propriétés analogues à celles de $[F_2]$ à un point de $(C_{1, 2})$.

\textbf{6.} \emph{Points fondamentaux} (6---7). Il y aura en général sur les deux surfaces des points dont les homologues restent indéterminés, de façon que tous les points d'une courbe correspondent à un de ces points. On les appelle les \emph{points fondamentaux} ou \emph{principaux}. Tout point dont on trouve plus d'un point correspondant doit être un point fondamental. A un point infiniment voisin d'un point fondamental correspond un point infiniment voisin de la courbe correspondante. Comme le point infiniment voisin détermine une direction par le point fondamental, les directions par ce point correspondent aux points de la courbe correspondante. A une courbe $(T_1)$ qui sur une des surfaces $[F_1]$ passe $\lambda$ fois par un de ses points fondamentaux $D_1$, correspond sur l'autre surface $[F_2]$ une courbe composée de la courbe $(D_{1, 2})$ qui correspond au point fondamental et d'une autre courbe $(T_{1, 2})$ qui rencontre $(D_{1, 2})$ aux $\lambda$ points correspondant aux directions des tangentes de $(T_1)$ au point $D_1$. Tout autre point de rencontre de $(D_{1, 2})$ et de $(T_{1, 2})$ sera --- à moins que l'intersection ne soit due seulement à une multiplicité de $[F_2]$ à ce point --- un point fondamental de $[F_2]$; car il correspond à la fois à $D_1$ et à un autre point de $[F_1]$.

Les tangentes d'une surface à un point fondamental forment un plan tangent si le point est simple, et un cône s'il est conique. Il faut dans le premier cas que la courbe correspondante de l'autre surface, dont les points correspondent aux directions des tangentes au point fondamental, soit du genre zéro, et dans le second, qu'elle soit du même genre que le cône tangent.

\textbf{7.} Afin de ne pas étendre trop nos recherches nous supposerons qu'il n'y ait sur aucune des deux surfaces \emph{aucun point fondamental doué d'une direction fondamentale}, c'est-à-dire un point où toutes les courbes dont les correspondantes rencontrent une courbe fixe sur l'autre surface, sont tangentes à une même direction. \emph{Cette restriction nous défendra d'attribuer aux cônes tangents aux points fondamentaux coniques d'autres génératrices doubles ou cuspidales que $1^0$ celles qui y sont tangentes à la courbe double ou cuspidale, et $2^0$ celles qui correspondent à des points doubles ou cuspidaux de la courbe correspondant au point fondamental}${}^{*)}$.

\textbf{*)} En effet, dans la direction d'une génératrice double du cône tangent à un point multiple qui n'est pas tangente à la courbe double, il n'y a qu'un seul point

\origpage{25}
Nous commencerons par \emph{supposer que tous les points coniques des surfaces soient fondamentaux}, et nous montrerons après que les résultats obtenus s'appliquent aussi à des cas où il y a sur les deux surfaces des points coniques qui se correspondent. De l'autre côté, nous appliquerons les notations $\mu$, $v$, $y$, $\eta$, $z$, $\xi$ non seulement aux points coniques, mais aussi aux points fondamentaux simples ($\mu = 1$, $v = \text{etc.} = 0$). Nous supposerons qu'il n'y ait \emph{aucun point fondamental qui se trouve à un point-pince ou à un point non-conique d'une courbe cuspidale}${}^{*)}$. Pour le moment nous supposerons encore que \emph{les surfaces ne soient pas douées de points-clos} ($\chi = 0$).

Nous désignerons par $\mu_{1, 2}$ l'ordre et par $\nu_{1, 2}$ la classe de la courbe $(D_{1, 2})$ qui sur $[F_2]$ correspond à un point fondamental $D_1$ de $[F_1]$; la courbe $(D_{1, 2})$ a, suivant les suppositions déjà faites, $\eta_1$ points doubles et $\xi_1$ points cuspidaux correspondant aux génératrices doubles et cuspidales du cône tangent au point $D_1$ qui ne sont pas tangentes à la courbe double et à la courbe cuspidale, et nous lui attribuerons encore $y_{1, 2}$ points doubles, apparents et propres, et $z_{1, 2}$ points cuspidaux${}^{**)}$. Alors on a, en désignant encore par $\pi_1$ le genre commun au cône tangent à $D_1$ et à la courbe $(D_{1, 2})$, les équations suivantes (où l'équation (V) de l'introduction du précédent article est renfermée)
\[
\left\{
\begin{aligned}
2(\pi_1 - 1) &= \nu_1 + z_1 + \xi_1 - 2\mu_1 = \mu_1(\mu_1 - 3) - 2(y_1 + \eta_1 + z_1 + \xi_1) \\
&= \nu_{1, 2} + z_{1, 2} + \xi_1 - 2\mu_{1, 2} = \mu_{1, 2}(\mu_{1, 2} - 3) - 2(y_{1, 2} + \eta_1 + z_{1, 2} + \xi_1).
\end{aligned}
\right. \tag{2}
\]

Les caractéristiques Plückeriennes de la courbe $(D_{2, 1})$ qui sur $[F_1]$ correspond à un point fondamental $D_2$ de $[F_2]$ seront désignées de la manière analogue; elles satisfont aux équations analogues aux équations (2).

\textbf{(* de p.~24)} consécutif. Par conséquent, si dans le cas où le point conique est fondamental, deux points distincts de sa courbe correspondante correspondaient aux génératrices du cône tangent qui coïncident dans la génératrice double, l'unique point consécutif dans cette direction aurait deux points correspondants et serait donc un nouveau point fondamental coïncidant avec le premier, et la direction deviendrait fondamentale. Le cas d'une génératrice cuspidale peut être regardée comme une limite de celui d'une génératrice double.

On pourrait, du reste, obtenir les mêmes résultats par une continuation des procédés analytiques de M.~Nöther. (Voir la première note du présent article.)

\textbf{*)} On peut regarder séparément les deux nappes qui passent par les points de la courbe double, tant que les points de la surface qui y coïncident ne sont pas consécutifs. (Aux points-pinces et aux points d'intersection avec la courbe cuspidale les points coïncidents sont consécutifs.) Un de deux points coïncidents, mais non pas consécutifs, peut être fondamental; s'ils le sont tous deux, on a seulement deux points fondamentaux simples.

\textbf{**)} Les points doubles propres de $(D_{1, 2})$ qui ne sont pas au nombre des $\eta_1$ que nous venons de nommer et les $z_{1, 2}$ points cuspidaux, se trouvent \emph{ou} aux points fondamentaux de $[F_2]$ \emph{ou} sur sa courbe double et sa courbe cuspidale.

\origpage{26}
\textbf{8.} A une courbe sur l'une des deux surfaces (dont aucune partie continue ne correspond à un point fondamental de l'autre surface et qui ne passe pas elle-même par un point fondamental), correspond une courbe qui est du même \emph{genre}, parce que les points des deux courbes se correspondent un-à-un.

A chaque intersection de deux courbes sur l'une des surfaces qui n'est pas due à une multiplicité de cette surface au point d'intersection et qui n'a pas lieu à un point fondamental, correspond sur l'autre surface une intersection des deux courbes correspondantes. Lorsque deux des points d'intersection situés sur l'une des surfaces coïncident, la même chose doit avoir lieu pour deux points d'intersection des deux courbes correspondantes. On en conclut que les points de contacts et les points doubles et cuspidaux des courbes de l'une des surfaces qui se trouvent à des points simples et non-fondamentaux, ont pour correspondants des points de contact et des points doubles et cuspidaux des courbes correspondantes sur l'autre surface.

\section*{II. Propriétés des courbes qui sur l'une des deux surfaces correspondent à des sections planes de l'autre.}

\textbf{9.} \emph{L'ordre} $s$ d'une courbe $(S_1)$ qui sur la surface $[F_1]$ correspond à une section plane de $[F_2]$, est égal au nombre des points où elle rencontre une section plane de $[F_1]$. $s$ est donc le nombre des points d'une section plane de l'une des surfaces dont les correspondants se trouvent sur une section plane de l'autre. On voit ainsi que \emph{l'ordre $s$ d'une courbe $(S_1)$ est égal à celui d'une courbe $(S_2)$ qui sur $[F_2]$ correspond à une section plane de $[F_1]$}.

\textbf{10.} Nous désignons par $r_{s, 1}$ la classe de la courbe $(S_1)$, par $n_{s, 1}$ le nombre des plans osculateurs à elle qui sortent d'un point quelconque, et par $h_{s, 1}$ le nombre des génératrices doubles d'un cône qui la projette d'un sommet quelconque. Les remarques préliminaires nous montrent \emph{qu'elle n'est pas, en général, douée de points cuspidaux}, et qu'elle passe $\mu_{1, 2}$ fois par un point fondamental, $\mu_{1, 2}$ ayant la signification indiquée au \no 7. Les points fondamentaux comptent donc dans le nombre $h_{s, 1}$ pour $\Sigma \frac{\mu_{1, 2}(\mu_{1, 2} - 1)}{2}$, où la somme $\Sigma$ est étendue à tous les points fondamentaux de $[F_1]$.

Comme le genre de $(S_1)$ est égal à celui d'une section plane de $[F_2]$, que nous avons appelé $p_2$, on trouve, en appliquant les formules de Plücker et de Clebsch à un cône qui projette $(S_1)$, les expressions suivantes:

\origpage{27}
\[
\begin{cases}
\begin{aligned}
r_{s, 1} &= 2(p_2 - 1) + 2s, \\
2h_{s, 1} &= s(s - 3) - 2(p_2 - 1), \\
n_{s, 1} &= 6(p_2 - 1) + 3s.
\end{aligned}
\end{cases} \tag{$3_a$}
\]
On trouve de même en appliquant les notations analogues à une courbe $(S_2)$ qui sur $[F_2]$ correspond à une section plane de $[F_1]$,
\[
\begin{cases}
\begin{aligned}
r_{s, 2} &= 2(p_1 - 1) + 2s, \\
2h_{s, 2} &= s(s - 3) - 2(p_1 - 1), \\
n_{s, 2} &= 6(p_1 - 1) + 3s.
\end{aligned}
\end{cases} \tag{$3_b$}
\]

\textbf{11.} Comme toute courbe $(S_1)$ passe $\mu_{1, 2}$ fois par un point fondamental de $[F_1]$, deux courbes $(S_1)$ se rencontrent $\Sigma(\mu_{1, 2}^2)$ fois aux points fondamentaux. Elles se rencontrent encore aux points correspondant aux $n_2$ intersections des deux sections planes de $[F_2]$ qui correspondent aux deux courbes $(S_1)$. Deux courbes $(S_1)$ se rencontrent donc, en général, en
\[
\Sigma(\mu_{1, 2}^2) + n_2
\]
points, et à ce nombre s'ajoute seulement en des cas particuliers des intersections „impropres“ à des points de la courbe double ou cuspidale de $[F_1]$${}^{*)}$.

On trouve de même que deux courbes $(S_2)$ ont, en général,
\[
\Sigma(\mu_{2, 1}^2) + n_1
\]
intersections.

\textbf{12.} Les deux cônes qui d'un même sommet $O$ projettent deux courbes $(S_1)$, étant de l'ordre $s$, auront $s^2$ génératrices communes. A ce nombre appartiennent les droites qui joignent $O$ aux $\Sigma(\mu_{1, 2}^2) + n_2$ points d'intersection des deux courbes $(S_1)$, et, dans le cas où les deux courbes sont infiniment voisines, on peut aussi trouver une expression du nombre de leurs \emph{intersections apparentes}, qu'il faut y ajouter pour avoir celui de toutes les génératrices communes. En égalant celui-ci à $s^2$, on trouve une équation à laquelle doivent satisfaire les nombres qui caractérisent les deux surfaces et les rapports de l'une à l'autre.

Le nombre des intersections apparentes de deux courbes $(S_1)$ infiniment voisines l'une de l'autre, sera composé $1^0$ de celui des
\[
\left[ h_{s, 1} - \Sigma\left( \frac{\mu_{1, 2}(\mu_{1, 2} - 1)}{2} \right) \right]
\]
points doubles apparents${}^{**)}$ de l'une de ces courbes pris \emph{deux fois}, et

\textbf{*)} On voit que les mots „intersections impropres“ et „intersections apparentes“ ont des significations différentes. Le nombre des intersections apparentes de deux courbes est celui des droites qui d'un même sommet projettent à la fois des points des deux courbes qui ne coïncident pas dans un point d'intersection.

\textbf{**)} Y compris les points doubles impropres s'il y a des courbes doubles qui se correspondent.

\origpage{28}
$2^0$ de celui des génératrices de contact d'un cône $[O\,(S_1)]$, projetant une courbe $(S_1)$, avec le cône enveloppe de la serie des cônes qui projettent toutes les courbes $(S_1)$ qui correspondent à un faisceau de sections planes de $[F_2]$. Ce cône enveloppe se compose $1^0$ du cône circonscrit à la surface $[F_1]$, et $2^0$ du cône projetant la courbe cuspidale $(C_1)$ de cette surface.

Pour trouver le nombre de contacts d'un cône $[O\,(S_1)]$ avec le cône circonscrit à $[F_1]$ mené du même sommet $O$, on fait usage de la première polaire de $O$ par rapport à $[F_1]$, qui est une surface de l'ordre $n_1 - 1$, passant par la courbe de contact $(A_1)$ du cône circonscrit, par la courbe double $(B_1)$ et par la courbe cuspidale $(C_1)$ de la surface $[F_1]$, tangente à elle le long de cette dernière courbe, et ayant pour $(\mu_1 - 1)$-tuple tout point $\mu_1$-tuple de la surface donnée $[F_1]$. Cette polaire rencontre la courbe $(S_1)$ à $s(n_1 - 1)$ points dont les $\Sigma[\mu_{1, 2}(\mu_1 - 1)]$ sont aux points fondamentaux. Les autres sont les points (non-fondamentaux) où $(S_1)$ rencontre les courbes $(A_1)$, $(B_1)$ et $(C_1)$. Les points d'intersection avec la courbe de contact $(A_1)$ déterminent les génératrices cherchées le long desquelles le cône $[O\,(S_1)]$ est tangent au cône circonscrit $[O\,(A_1)]$. Les points d'intersection avec la courbe $(B_1)$ correspondent aux points d'intersection d'une section plane de la surface $[F_2]$ avec la courbe $(B_{1, 2})$${}^{*)}$ correspondante à $(B_1)$, et leur nombre est donc égal à l'ordre $b_{1, 2}$ de $(B_{1, 2})$. De même $(S_1)$ rencontre${}^{**)}$ la courbe cuspidale $(C_1)$ à $c_{1, 2}$ points, $c_{1, 2}$ étant l'ordre de la courbe $(C_{1, 2})$ correspondant à $C_1$; mais à chacun de ces $c_{1, 2}$ points coïncident (suivant le \no 4) trois points d'intersection de $(S_1)$ avec la polaire. On trouve ainsi que le cône $[O\,(S_1)]$ qui projette $(S_1)$ a
\[
s(n_1 - 1) - \Sigma [\mu_{1, 2}(\mu_1 - 1)] - b_{1, 2} - 3c_{1, 2}
\]
contacts avec le cône circonscrit $[O\,(A_1)]$.

Les génératrices de contact du cône $[O\,(S_1)]$ avec le cône $[O\,(C_1)]$ se déterminent au moyen des $c_{1, 2}$ points de contact de la courbe $(S_1)$ avec la courbe $(C_1)$, et leur nombre est donc $c_{1, 2}$.

On obtient ainsi l'équation suivante
\[
\begin{aligned}
s^2 &= \Sigma(\mu_{1, 2}^2) + n_2 + 2\left[ h_{s, 1} - \Sigma\left( \frac{\mu_{1, 2}(\mu_{1, 2} - 1)}{2} \right) \right] \\
&\quad + \{s(n_1 - 1) - \Sigma[\mu_{1, 2}(\mu_1 - 1)] - b_{1, 2} - 3c_{1, 2}\} + c_{1, 2}.
\end{aligned}
\]
En y substituant l'expression (3) de $h_{s, 1}$ on trouve l'équation suivante

\textbf{*)} Si $(B_1)$ ou $(C_1)$ passe par des points fondamentaux, nous ne regardons pas comme faisant partie de $(B_{1, 2})$ ou de $(C_{1, 2})$ les courbes qui correspondent à ces points fondamentaux.

\textbf{**)} Elle y est tangente; voir le \no 3.

\origpage{29}
\[
n_2 - 2(p_2 - 1) + s(n_1 - 4) - \Sigma [\mu_{1, 2}(\mu_1 - 2)] - b_{1, 2} - 2c_{1, 2} = 0. \tag{$\mathrm{I}_a^{*)}$}
\]
On trouve, en remplaçant les deux surfaces $[F_1]$ et $[F_2]$ l'une par l'autre
\[
n_1 - 2(p_1 - 1) + s(n_2 - 4) - \Sigma [\mu_{2, 1}(\mu_2 - 2)] - b_{2, 1} - 2c_{2, 1} = 0. \tag{$\mathrm{I}_b$}
\]

\section*{III. Construction et propriétés de surfaces auxiliaires.}

\textbf{13.} \emph{Construction des surfaces auxiliaires.} Dans nos recherches ultérieures nous ferons usage d'une surface auxiliaire $[F_3]$ définie de la manière suivante: \emph{les droites qui joignent les points de $[F_1]$ à un point fixe $O_3$ rencontrent les plans qui joignent les points correspondants de $[F_2]$ à une droite fixe $(M_3)$, aux points de la surface $[F_3]$}. La surface $[F_3]$ sera donc engendrée par le faisceau de plans passant par $(M_3)$, et par la série de cônes $[O_3\,(S_1)]$ qui joignent $O_3$ aux courbes $(S_1)$ qui sur $[F_1]$ correspondent aux sections faites à $[F_2]$ par les plans du faisceau: chacun des cônes correspond à un plan du faisceau qu'il rencontre dans une courbe située sur $[F_3]$.

Nous ferons de même usage d'une surface $[F_4]$ qu'on construit de la même manière que $[F_3]$ en substituant seulement les deux surfaces l'une à l'autre. Nous désignerons par $O_4$ et $(M_4)$ le point et le plan fixes dont on fait usage dans la construction de $[F_4]$. Là, où nous n'indiquerons pas directement les résultats qui ont égard à $[F_4]$, on les déduit sans difficulté de ceux qui ont égard à $[F_3]$.

Suivant la construction des surfaces $[F_3]$ et $[F_4]$, \emph{les points de toutes les quatre surfaces $[F_1]$, $[F_2]$, $[F_3]$ et $[F_4]$ se correspondent un-à-un}.

\textbf{14.} \emph{Ordre de la surface $[F_3]$.} Désignons par $P$ et $Q$ les points où une droite fixe $(L)$ rencontre une droite par $O_3$ et un plan par $(M_3)$ qui sont déterminés par des points correspondants de $[F_1]$ et de $[F_2]$. Alors, comme une droite $(A\,P)$ rencontre $[F_1]$ à $n_1$ points, à un point $P$ correspondent $n_1$ points $Q$, et, comme un plan $[(M_3)\,Q]$ contient $s$ points de $[F_2]$ dont les correspondants se trouvent dans le plan $[O_3\,(L)]$, à un point $Q$ correspondent $s$ points $P$. Il y aura donc, suivant le principe de correspondance, sur $(L)$ $n_1 + s$ coïncidences de points correspondants $P$ et $Q$, et le nombre des intersections de $(L)$ avec $[F_3]$, ou bien l'ordre de cette surface, sera par conséquent
\[
n_3 = n_1 + s.
\]

On aurait pu se contenter de chercher le nombre des intersections d'une droite passant par $O_3$ avec la surface $[F_3]$. Une droite par $O_3$ rencontre $[F_3]$ $1^0$ aux $n_1$ points qui sur cette nouvelle surface corres-

\textbf{*)} Si la courbe $(B_1)$ avait été $k$-tuple, on aurait dû donner au terme $- b_{1, 2}$ le coefficient $(k - 1)$.

\origpage{30}
pondent aux intersections de la même droite avec la surface $[F_1]$, et $2^0$ aux $s$ points coïncidant avec $O_3$ où elle rencontre le cône $[O_3\,(S_1)]$ déterminé par la courbe $(S_1)$ qui correspond à la section faite à $[F_2]$ par le plan $[(M_3)\,O_3]$. On voit ainsi en même temps que \emph{la surface $[F_3]$ a, au point $O_3$, un point $s$-tuple}.

On aurait pu de même se contenter de chercher l'ordre de la section faite à $[F_3]$ par un plan passant par la \emph{droite} $(M_3)$. Cette section est composée $1^0$ de la courbe d'ordre $s$ où son plan rencontre le cône qui y correspond dans la série de cônes $[O_3\,(S_1)]$, et $2^0$ d'une droite multiple coïncidant avec $(M_3)$. Celle-ci est \emph{$n_1$-tuple}; car à un de ses points $R$ la surface $[F_3]$ a pour plans tangents ceux qui joignent $(M_3)$ aux $n_1$ points de $[F_2]$ qui correspondent aux $n_1$ points de rencontre de la droite $(O_3\,R)$ avec $[F_1]$.

La surface $[F_4]$ est de l'ordre $n_4 = n_2 + s$; le point $O_4$ y est $s$-tuple et la droite $(M_4)$ $n_2$-tuple.

\textbf{15.} \emph{Le point singulier $O_3$ de la surface $[F_3]$.} (15---17). Nous avons vu que le point $O_3$ est $s$-tuple sur la surface $[F_3]$, et la démonstration elle-même nous montre que le cône tangent à ce point est le cône $[O_3\,(S_1)]$ qui joint $O_3$ à la courbe $(S_1)$ qui sur $[F_1]$ correspond à la section faite à $[F_2]$ par le plan $[(M_3)\,O_3]$. Nous avons déjà (dans le \no 10) parlé des caractéristiques Plückeriennes de ce cône tangent. Nous avons encore à faire remarquer que les génératrices doubles de ce cône sont tangentes aux branches de la courbe double de $[F_3]$ qui passent par le point $O_3$, ce qu'on voit sans difficulté, et que \emph{les génératrices suivantes se trouvent en entier sur la surface $[F_3]$: $1^0$ celles qui passent par les points fondamentaux de $[F_1]$, chacune prise $\mu_{1, 2}$ fois, $2^0$ celles qui passent par les $n_2$ points de $[F_1]$ qui correspondent aux points où $[F_2]$ est rencontrée par la droite $(M_3)$, et $3^0$ celles qui passent par les $s$ points de la section que fait le plan $[O_3\,(M_3)]$ à la surface $[F_1]$ dont les correspondants sur $[F_2]$ se trouvent dans le même plan}.

En effet, les points de la surface $[F_1]$ qui déterminent les deux premières classes de droites, se trouvent, comme nous l'avons prouvé aux \nos 10 et 11 --- $\mu_{1, 2}$ ou une fois --- sur toutes les courbes $(S_1)$ qui correspondent aux sections faites à $[F_2]$ par les plans passant par $(M_3)$, et les droites qui y joignent le point $O_3$ se trouvent par conséquent --- le même nombre de fois --- sur tous les cônes $[O_3\,(S_1)]$ de la série. Quant à la dernière classe de droites elles se trouvent à la fois dans le plan $[(M_3)\,O_3]$ et sur le cône qui y correspond.

\textbf{16.} Chacun des points de $[F_1]$ qui déterminent ces droites devient un point fondamental dans la correspondance de $[F_1]$ avec la surface $[F_3]$, et les directions par les points sur $[F_1]$ correspondent aux points des droites sur $[F_3]$. Le plan par une des droites sur $[F_3]$ déterminé

\origpage{31}
par une direction par le point correspondant de $[F_1]$, est tangent à $[F_3]$ au point qui correspond à la direction. On a ainsi le moyen de déterminer les nombres caractéristiques${}^{*)}$ de ces droites situés sur la surface $[F_3]$.

Regardons premièrement la droite $(D_{1, 3})$ qui joint $O_3$ à un point $(D_1)$ qui est fondamental sur $[F_1]$ déjà dans sa correspondance avec l'autre surface donnée $[F_2]$. Nous avons déjà indiqué que $(D_{1, 3})$ est une droite $\mu_{1, 2}$-tuple de $[F_3]$, et l'on trouve les $\mu_{1, 2}$ plans tangents de $[F_3]$ à un point quelconque $P$ de $(D_{1, 3})$ au moyen du plan $[(M_3)\,P]$ qui joint le point $P$ à la droite fixe $(M_3)$: ce plan rencontre la courbe $(D_{1, 2})$ qui sur $[F_2]$ correspond à $D_1$ à $\mu_{1, 2}$ points dont les directions correspondantes par le point $D_1$ déterminent les plans tangents cherchés. Réciproquement, un plan passant par la droite $(D_{1, 3})$ contient $\mu_1$ directions par le point fondamental: il sera donc tangent à la surface $[F_3]$ aux $\mu_1$ points que déterminent les plans joignant $(M_3)$ aux points de la courbe $(D_{1, 2})$ qui correspondent à ces directions. Deux des plans tangents à $[F_3]$ à un même point $P$ de $(D_{1, 3})$ sont consécutifs dans le cas où le plan joignant $P$ à $(M_1)$ est tangent à la courbe $(D_{1, 2})$, ou passe par un point cuspidal de cette courbe, ou par un des $\eta_1$ points doubles${}^{**)}$ dont les génératrices correspondantes du cône tangent à $D_1$ sont doubles sans être tangentes à la courbe double de $[F_1]$; et deux points de contact d'un même plan par $(D_{1, 3})$ sont consécutifs dans le cas où le plan est tangent au cône tangent au point $D_1$, ou passe par une des génératrices cuspidales de ce cône, ou par une des $\eta_1$ génératrices doubles qui ne sont pas tangentes à la courbe double de $[F_1]$. On trouve ainsi --- en observant que, la surface $[F_3]$ étant construite comme lieu de points, les coïncidences auront lieu des manières qui sont les plus générales pour des surfaces regardées comme lieux de points${}^{***)}$ --- les valeurs suivantes des caractéristiques de $(D_{1, 3})$:

\textbf{*)} Nous renvoyons, sur la signification de ces caractéristiques et sur les notations que nous y appliquerons aussi ici, à notre précédent article, notamment au résumé qui se trouve à son \no 10.

\textbf{**)} Les notations qui ont égard aux points fondamentaux de $[F_1]$ sont indiquées au \no 7 du présent article.

\textbf{***)} Voir dans le précédent article les \nos 3, 4 et 5. On pourrait croire que les points de $(D_{1, 3})$, déterminés par les $\eta_1$ points doubles et les $\xi_1$ points cuspidaux de $(D_{1, 2})$ dont les génératrices correspondantes du cône tangent au point $D_1$ sont doubles et cuspidales sans être tangentes à la courbe double et à la courbe cuspidale de $[F_1]$, appartenaient aux nombres $\bar{f}$ et $\bar{i}$ des points d'intersection de $(D_{1, 3})$ avec la courbe double et avec la courbe cuspidale de $[F_3]$ (voir le \no 6 du précédent article). On voit que cela n'a pas lieu, en regardant la section de $[F_3]$ que fait le plan joignant $(M_3)$ à un de ces $(\eta_1 + \xi_1)$ points. Celle-ci aura à ce point un point cuspidal, et non pas un point de contact de deux

\origpage{32}
\[
\bar{n} = \mu_{1, 2}, \quad \bar{n}' = \mu_1, \quad 1\bar{\jmath} = \nu_{1, 2} + z_{1, 2}, \quad 2\bar{\jmath} = \eta_1, \quad 3\bar{\jmath} = \xi_1.
\]
\[
1\bar{\beta}' = \nu_1 + z_1, \quad \bar{\jmath}' = \bar{\beta} = \bar{i}' = 0.
\]
En cherchant aussi les points de $(D_{1, 3})$ et les plans par $(D_{1, 3})$ dont deux plans tangents ou deux points de contact coïncident sans être consécutifs, on trouverait le nombre $\bar{f}$ de contacts de deux des nappes de $[F_3]$ qui passent par $(D_{1, 3})$. Ce nombre se trouve aussi au moyen de l'une des équations (5) du précédent article; l'autre de ces équations donnera
\[
2\mu_{1, 2} + \nu_1 + z_1 = 2\mu_1 + \nu_{1, 2} + z_{1, 2},
\]
qui est une des équations (2) du présent article. Les équations (7) et (8) du précédent article, qui subiront des altérations à cause du point singulier $O_3$ qui se trouve sur la droite $(D_{1, 3})$, ne nous donneront à présent aucun résultat nouveau.

Des procédés analogues nous montrent que les autres $n_2 + s$ droites par $O_3$ qui se trouvent en entier sur $[F_3]$ sont simples, et qu'un plan passant par une de ces droites n'a, en général, qu'un seul contact avec $[F_3]$. On a donc pour chacune de ces droites $\bar{n} = \bar{n}' = 1$, $\bar{\jmath} = \bar{\jmath}' = \bar{\beta} = \bar{\beta}' = \bar{i} = \bar{f} = 0$.

\textbf{17.} \emph{Le cône tangent d'une surface à un point $s$-tuple} est le lieu des droites par le point qui contiennent $s + 1$ points coïncidents de la surface. Il y a $s(s + 1)$ génératrices du cône qui contiennent $s + 2$ points coïncidents. Pour le point $O_3$, $s$-tuple sur la surface $[F_3]$, ces $s(s + 1)$ génératrices seront: $1^0$ les $s^2$, dont nous avons déjà rendu compte au \no 12, où le cône correspondant à la section faite à $[F_2]$ par le plan $[(M_3)\,O_3]$ rencontre le cône consécutif de la série de cônes $[O_3\,(S_1)]$ qui correspondent aux plans par $(M_3)$, et $2^0$ les $s$ droites d'intersection du même cône avec le plan correspondant $[(M_3)\,O_3]$.

\textbf{18.} \emph{La droite multiple $(M_3)$ de la surface $[F_3]$.} Nous avons vu que la droite $(M_3)$ est, sur la surface $[F_3]$, multiple de l'ordre $n_1$. Les plans tangents de la surface à un point quelconque $R$ de $(M_3)$ sont ceux qui passent par les points de la surface $[F_2]$ qui correspondent aux $n_1$ points de rencontre de la droite $(O_3\,R)$ avec la surface $[F_1]$. Les points de contact d'un plan quelconque $[Q]$ par $(M_3)$, sont les points de rencontre de $(M_3)$ avec les droites qui joignent $O_3$ aux $s$ points de la courbe d'intersection du plan $[O_3\,(M_3)]$ avec $[F_1]$ dont les correspondants sur $[F_2]$ se trouvent dans le plan $[Q]$. Ces derniers points sont les points d'intersection du plan $[Q]$ avec la courbe $(S_2)$

\textbf{(* de p.~31)} branches ou un point de rebroussement de seconde espèce. --- Nous verrons, du reste, que $[F_3]$ n'a aucune courbe cuspidale, de façon que $\bar{i}$ (ainsi que $\bar{\beta}$) doive être égal à zéro.

\origpage{33}
qui sur $[F_2]$ correspond à la section faite à $[F_1]$ par le plan $[O_3\,(M_3)]$. En cherchant encore les points de $(M_3)$ où deux plans tangents sont consécutifs, et les plans par $(M_3)$ dont deux points de contact sont consécutifs --- ce qui ne causera aucune difficulté --- on trouve pour la droite multiple $(M_3)$ les nombres caractéristiques suivants:
\[
\bar{n} = n_1, \quad \bar{n}' = s, \quad \bar{\jmath} = a_1 + c_1, \quad \bar{\beta}' = r_{s, 2},
\]
\[
2\bar{\jmath} = 3\bar{\jmath} = \bar{\jmath}' = \bar{\beta} = \bar{i} = 0.{}^{*)}
\]

Le nombre $\bar{f}$ indiquera l'ordre d'un cône, lieu de droites joignant deux points de $[F_1]$ dont les correspondants sur $[F_2]$ se trouvent sur des droites qui rencontrent une droite fixe. Le sommet du cône est ici à $O_3$, et la droite fixe est à $(M_3)$.

Les formules (5) du précédent article nous donneront ici, $1^0$ qu'une courbe $(S_2)$ est du même genre qu'une section plane de $[F_1]$, ce que nous avons déjà fait remarquer, et $2^0$ l'expression suivante de $\bar{f}$:
\[
2\bar{f} = 2\,n_1 s - 2\,s - a_1 - c_1,
\]
ou bien (voir les formules (1) du présent article)
\[
\bar{f} = (n_1 - 1)(s - 1) - p_1.
\]

Le même résultat se trouve sans difficulté au moyen du principe de correspondance. --- L'équation (7) du précédent article ne nous donnera que le résultat évident que le nombre $\bar{t}$ des intersections de la droite $(M_1)$ avec d'autres nappes de la surface que les $n_1$ qui passent par elle, est égal à zéro. L'équation (8) du précédent article nous sera utile dans ce qui suit.

\textbf{19.} \emph{Section plane de la surface $[F_3]$.} On trouve les caractéristiques Plückeriennes d'une section plane de la surface $[F_3]$ au moyen du précédent article${}^{**)}$, où nous avons exprimé les caractéristiques Plückeriennes de la section \emph{résidue} d'un plan passant par la droite multiple --- c'est-à-dire de ce qui reste de la courbe d'intersection lorsqu'on fait abstraction de la droite multiple --- au moyen de celles d'une section quelconque. Nous ferons ici, des mêmes expressions, l'usage inverse, en sachant que la section résidue faite à la surface $[F_3]$ par un plan passant par la droite multiple $(M_3)$ est la courbe d'intersection de ce plan avec un cône $[O_3\,(S_1)]$, et qu'elle est par conséquent de l'ordre $s$, de la classe $r_{s, 1}$ et dénuée de points cuspidaux. Ce procédé est iden-

\textbf{*)} Dans le cas, que nous avons negligé, où il y a sur les deux surfaces des courbes cuspidales correspondant l'une à l'autre, $\bar{i}$ sera l'ordre de celle de la surface $[F_1]$.

\textbf{**)} Voir le résumé à son \no 10.

\origpage{34}
tique au troisième de ceux qui nous ont donné l'ordre $n_3 = n_1 + s$ de la surface $[F_3]$. On trouve de même, en désignant par $a_3$ la classe d'une section plane de $[F_3]$,
\[
a_3 = r_{s, 1} + 2\,s + a_1 + c_1
\]
ou bien, en réduisant cette expression au moyen des formules (1) et (3) du présent article
\[
\begin{aligned}
a_3 &= 4\,s + 2\,n_1 + 2(p_1 - 1) + 2(p_2 - 1) \\
&= a_1 + c_1 + a_2 + c_2 - 2\,n_2 + 4\,s.
\end{aligned} \tag{4}
\]
On trouve enfin
\[
c_3 = 0, \tag{5}
\]
de façon que $[F_3]$ soit dénuée de courbe cuspidale. Les autres caractéristiques Plückeriennes de la section dépendent de $n_3$, de $a_3$ et de $c_3$.

On aurait trouvé les mêmes résultats en regardant la section faite à $[F_3]$ par un plan passant par son point multiple $O_3$.

\textbf{20.} On voit sans difficulté que le point $O_3$ est le seul point fondamental de $[F_3]$ dans sa correspondance avec $[F_1]$ ou $[F_2]$, et qu'il a pour courbes correspondantes, sur $[F_2]$ la section faite par le plan $[O_3\,(M_3)]$, et sur $[F_1]$ la courbe $(S_1)$ qui correspond à cette section et qui détermine le cône tangent de $[F_3]$ au point $O_3$. Nous avons indiqué au \no 16., que, dans la correspondance de $[F_3]$ avec $[F_1]$, les points fondamentaux de $[F_1]$ sont ceux qui déterminent les droites par $O_3$ qui se trouvent en entier sur $[F_3]$, et que les courbes correspondantes sont ces droites. On en déduit que les points fondamentaux de $[F_2]$ dans sa correspondance avec $[F_3]$ sont, $1^0$ ceux qui le sont dans sa correspondance avec $[F_1]$, $2^0$ les $n_2$ points où elle est rencontrée par la droite $(M_3)$, et $3^0$ les $s$ points de la section faite par le plan $[O_3\,(M_3)]$, dont les correspondants sur $[F_1]$ se trouvent dans le même plan.

L'application de la formule (I) (au \no 12) à la correspondance de la surface $[F_3]$ avec $[F_1]$ ou $[F_2]$ ne donne aucun résultat nouveau et important, mais elle peut servir de vérification des résultats déjà trouvés.

\section*{IV. Cônes circonscrits à $[F_3]$ ou à $[F_4]$; usage des surfaces auxiliaires.}

\textbf{21.} Nous avons déjà trouvé au \no 19. l'ordre $a_3$ d'un cône circonscrit à $[F_3]$, qui est égal à la classe d'une section plane. Deux autres caractéristiques Plückeriennes du cône circonscrit seront déterminées, chacune par deux voies différentes, qui donneront deux expressions différentes. \emph{En égalant les deux expressions d'un même nombre on trouve les équations relatives aux deux surfaces données qui sont le but de la construction des surfaces auxiliaires}.

\origpage{35}
\textbf{22.} \emph{Classe de la surface $[F_3]$} (22---24). On trouve la classe $n'_3$ de la surface $[F_3]$ en appliquant la formule (8) du précédent article
\[
n' = \bar{n}' + \bar{n} + \bar{t}' + 2\cdot \bar{\beta}'
\]
à la droite multiple $(M_3)$ de la surface $[F_3]$. Alors $n'$ sera la classe cherchée $n'_3$; $\bar{n}'$, $\bar{n}$ et $\bar{\beta}'$ auront, comme nous l'avons indiqué au \no 18., les valeurs $s$, $n_1$ et $r_{s, 2}$. Il reste donc à chercher le nombre $\bar{t}'$ des plans $[Q]$ par la droite $(M_3)$ qui ont un $(s + 1)$-ième contact avec la surface $[F_3]$, ou des plans dont les sections résidues ont des points doubles qui ne sont ni au point multiple $O_3$ ni sur la courbe double de $[F_3]$. Or, la section résidue d'un plan $[Q]$ est la section que fait le cône de la série $[O_3\,(S_1)]$ qui correspond au plan $[Q]$, et ce plan sera donc un des $\bar{t}'$ plans cherchés si le cône correspondant a une génératrice double qui n'est pas au nombre des $h_{s, 1}$ dont tous les cônes $[O_3\,(S_1)]$ sont doués (voir le \no 10.), et qui déterminent seulement des points de la courbe double de $[F_3]$. Cela n'arrivera que lorsque la courbe $(S_1)$ obtient un point double propre autre que ceux qui sont aux points fondamentaux de $[F_1]$, ce qui a lieu dans les cas suivants: $1^0$ lorsque le plan $[Q]$ est tangent à $[F_2]$, $2^0$ lorsqu'il est tangent à la courbe cuspidale de $[F_2]$, et $3^0$ lorsqu'il passe par un point fondamental $D_2$ de $[F_2]$. En effet, dans les deux premiers cas, le point de $[F_1]$ qui correspond au point de contact sera un point double de $(S_1)$, suivant les \nos 8. et 5. Dans le troisième, la courbe $(S_1)$ se divisera, suivant le \no 6., en deux courbes qui se rencontrent aux $\mu_2$ points de $[F_1]$ correspondant aux $\mu_2$ directions du point fondamental (et $\mu_2$-tuple) $D_2$ qui se trouvent dans le plan $[Q]$, et la courbe $(D_{2, 1})$ qui correspond au point $D_2$ aura encore $\eta_2$ points doubles et $\xi_2$ points cuspidaux, correspondant aux génératrices doubles et cuspidales du cône tangent à $D_2$ qui ne sont pas tangentes à la courbe double et à la courbe cuspidale de $[F_2]$. On voit ainsi que dans le plan $[(M_3)\,D_2]$ coïncident $\mu_2 + \eta_2$ plans tangents simples et $\xi_2$ plans tangents stationnaires de $[F_3]$, de façon qu'il compte pour $\mu_2 + \eta_2 + 2\,\xi_2$ plans tangents simples. On aura donc${}^{*)}$:
\[
\bar{t}' = n'_2 + r_2 + \Sigma(\mu_2 + \eta_2 + 2\,\xi_2),
\]
et la formule du précédent article que nous venons de citer, nous donne ainsi
\[
n'_3 = s + n_1 + n'_2 + r_2 + \Sigma(\mu_2 + \eta_2 + 2\,\xi_2) + 2\,r_{s, 2} \tag{$6_a$}
\]
ou, suivant la première formule $(3_b)$ au \no 10.,
\[
n'_3 = n_1 + 4(p_1 - 1) + n'_2 + r_2 + \Sigma(\mu_2 + \eta_2 + 2\,\xi_2) + 5\,s \tag{$6_b$}
\]

\textbf{*)} On aurait trouvé le même résultat en regardant immédiatement la correspondance des surfaces $[F_2]$ et $[F_3]$ et des sections qu'y fait un même plan $[Q]$, au lieu de celle de $[F_2]$ avec $[F_1]$ et de la section plane de $[F_2]$ avec une courbe $(S_1)$.

\origpage{36}
ou bien, suivant la dernière formule (1) au \no 1.,
\[
n'_3 = 2\,a_1 + 2\,c_1 - 3\,n_1 + n'_2 + r_2 + \Sigma(\mu_2 + \eta_2 + 2\,\xi_2) + 5\,s. \tag{$6_c$}
\]

\textbf{23.} On obtient par une voie différente une autre expression de $n'_3$, qui est la classe, non seulement de la surface $[F_3]$, mais aussi d'un cône circonscrit mené à elle d'un sommet quelconque $E$. Ce cône aura la droite $(E\,O_3)$ pour génératrice multiple, le degré de multiplicité étant égal au nombre $r_{s, 1}$ des plans tangents au cône tangent au point $O_3$ qui passent par $(E\,O_3)$. On trouve donc la classe $n'_3$ du cône circonscrit en ajoutant le nombre $2\,r_{s, 1}$ à celui des plans qui passent par $(E\,O_3)$ et qui sont tangents à $[F_3]$ à des points différents de $O_3$. Nous allons chercher ce dernier nombre.

Comme les points correspondants des surfaces $[F_1]$ et $[F_3]$ se trouvent sur les mêmes droites par $O_3$, ils sont aussi sur les mêmes plans par $(E\,O_3)$, et les sections que fait un même plan $[R]$ par $(E\,O_3)$ sont, par conséquent, des courbes correspondantes. On voit donc, en y appliquant les mêmes raisonnements qu'au \no 22., que les plans par $(E\,O_3)$ tangents à $[F_3]$ à des points différents de $O_3$ sont: $1^0$ les $n'_1$ qui sont tangents à $[F_1]$, $2^0$ les $r_1$ qui sont tangents à la courbe cuspidale $(C_1)$ de cette surface, et $3^0$ ceux qui passent par les points de $[F_1]$ qui sont fondamentaux dans sa correspondance avec $[F_3]$. Ces derniers plans tangents à $[F_3]$ sont (voir le \no 20.) les plans passant par les droites de $[F_3]$ qui sortent de $O_3$. Si l'on emploie, pour déterminer leurs degrés de multiplicité, les mêmes règles${}^{*)}$ que dans le \no 22., on trouve qu'ils comptent pour $\Sigma(\mu_1 + \eta_1 + 2\,\xi_1) + n_2 + s$ plans tangents simples. Le nombre cherché des plans tangents à $[F_3]$ à des points différents de $O_3$, est donc
\[
n'_1 + r_1 + \Sigma(\mu_1 + \eta_1 + 2\,\xi_1) + n_2 + s.
\]
On trouve ainsi
\[
n'_3 = s + n_2 + n'_1 + r_1 + \Sigma(\mu_1 + \eta_1 + 2\,\xi_1) + 2\,r_{s, 1}. \tag{7}
\]

\emph{Autrement.} On pourrait trouver le même résultat d'une manière un peu différente. Les plans tangents à $[F_3]$ menés par $(E\,O_3)$ dont les points de contact ne sont pas à $O_3$, sont $1^0$ ceux qui sont tangents au cône lieu des droites tangentes à $(F_3)$ à des points différents de $O_3$, et $2^0$ ceux qui passent par les droites par $O_3$ qui se trouvent en entier sur $[F_3]$. Comme $[F_3]$ n'a aucune courbe cuspidale (voir au \no 19.), toute droite qui rencontre cette surface en deux points consécutifs y sera tangente, et dans cette condition seront les droites par $O_3$ qui ren-

\textbf{*)} La légitimité de l'application de ces règles au cas actuel où la courbe $(D_{1, 3})$ qui sur $[F_3]$ correspond à un point fondamental $D_1$ de $[F_1]$ dégénère en $\mu_{1, 2}$ droites coïncidentes, pourrait sembler douteuse; mais le résultat sera assuré par notre seconde démonstration.

\origpage{37}
contrent $[F_1]$ en deux points consécutifs. Le cône lieu de ces tangentes sera donc composé du cône $[O_3\,(A_1)]$ circonscrit à $[F_1]$ et du cône $[O_3\,(C_1)]$ qui projette sa courbe cuspidale. Il sera par conséquent de la classe $n'_1 + r_1$. Les plans passant par les droites de $[F_3]$ auront suivant le précédent article $\bar{n}' + 2\bar{\jmath} + 2\cdot 3\bar{\jmath}$ contacts avec cette surface, $\bar{n}'$, $2\bar{\jmath}$ et $3\bar{\jmath}$ prenant pour les droites respectives les valeurs indiquées au \no 16. Les plans joignant $(E\,O_3)$ aux droites de $[F_3]$ comptent donc pour $\Sigma(\mu_1 + \eta_1 + 2\,\xi_1) + n_2 + s$ plans tangents. ---

L'expression (7) se transformera par les mêmes substitutions que nous avons appliquées à l'expression (6) dans la suivante
\[
n'_3 = 2\,a_2 + 2\,c_2 - 3\,n_2 + n'_1 + r_1 + \Sigma(\mu_1 + \eta_1 + 2\,\xi_1) + 5\,s. \tag{$7_c$}
\]

\textbf{24.} En égalant les deux expressions ($6_c$) et ($7_c$) trouvées dans les \nos précédents on obtient l'équation suivante:
\[
\begin{aligned}
n'_1 - 2\,a_1 + r_1 - 2\,c_1 + 3\,n_1 + \Sigma(\mu_1 + \eta_1 + 2\,\xi_1) \\
= n'_2 - 2\,a_2 + r_2 - 2\,c_2 + 3\,n_2 + \Sigma(\mu_2 + \eta_2 + 2\,\xi_2).
\end{aligned} \tag{II}
\]
Cette relation est la même que j'ai exposée dans les \emph{Comptes Rendus de l'Académie des Sciences}${}^{*)}$ (T.~LXX, p.~745).

\textbf{25.} \emph{Identité des caractéristiques Plückeriennes de certains cônes circonscrits aux deux surfaces auxiliaires $[F_3]$ et $[F_4]$.} On voit, au moyen du précédent article${}^{**)}$ et des formules (4), ($6_c$) et ($7_c$) du présent article, que, pour un sommet $P$ situé sur la droite $(M_3)$, le cône \emph{résidu} circonscrit à $[F_3]$ est de l'ordre
\[
a_3 - 2\,n_1 = 4\,s + 2(p_1 - 1) + 2(p_2 - 1),
\]
et de la classe
\[
\begin{aligned}
n'_3 - s &= 2\,a_1 + 2\,c_1 - 3\,n_1 + n'_2 + r_2 + \Sigma(\mu_2 + \eta_2 + 2\,\xi_2) + 4\,s \\
&= 2\,a_2 + 2\,c_2 - 3\,n_2 + n'_1 + r_1 + \Sigma(\mu_1 + \eta_1 + 2\,\xi_1) + 4\,s.
\end{aligned}
\]
L'expression de l'ordre du cône est symétrique par rapport aux nombres caractéristiques des deux surfaces données $[F_1]$ et $[F_2]$, et les deux expressions de la classe se transforment l'une en l'autre par la substitution de l'une des surfaces $[F_1]$ et $[F_2]$ à l'autre. Les autres

\textbf{*)} Nous n'avons ici démontré la formule que dans les suppositions que nous avons déjà indiquées, au nombre desquelles se trouve celle qu'il n'y ait pas sur les surfaces données des courbes cuspidales correspondant l'une à l'autre. S'il y en a je crois avoir trouvé qu'on doit ajouter aux deux membres de l'équation les nombres respectifs des points cuspidaux des courbes cuspidales qui se correspondent sur les deux surfaces $[F_1]$ et $[F_2]$. J'aurais donc dû, dans les Comptes Rendus où je n'ai pas fait cette même restriction, trouver ce résultat un peu modifié. --- Dans les Comptes Rendus je n'attribuais aux surfaces que des points fondamentaux simples.

\textbf{**)} Voir le tableau à son \no 10.

\origpage{38}
caractéristiques Plückeriennes du même cône offriront la même symétrie, ce qu'il suffira --- à cause des relations Plückeriennes --- de prouver pour une seule de celles qui restent. Nous le ferons pour le nombre de ses arêtes cuspidales, c'est-à-dire pour le nombre des droites par le sommet $P$ qui ont un contact triponctuel avec $[F_3]$.

Soit $(L)$ une de ces droites. Alors il faut que le plan $[O_3\,(L)]$ contienne trois génératrices consécutives du cône $[O_3\,(S_1)]$ mené de $O_3$ à la courbe $(S_1)$ qui sur $[F_1]$ correspond à la section que fait le plan $[(M_3)\,(L)]$ à la surface $[F_2]$, ou bien, qu'il ait un contact triponctuel avec cette courbe $(S_1)$. Le nombre cherché est donc celui des plans d'un faisceau qui ont un contact triponctuel avec une courbe $(S_1)$ qui sur $[F_1]$ correspond à une section faite à $[F_2]$ par un plan d'un autre faisceau. Or le contact triponctuel d'un plan avec une courbe $(S_1)$ ou, ce qui est la même chose, celui de la section que fait le plan à $[F_1]$ avec $(S_1)$, amène aussi celui de la courbe $(S_2)$ qui sur $[F_2]$ correspond à la section plane de $[F_1]$ avec la courbe plane qui sur $[F_2]$ correspond à $(S_1)$. Le nombre des arêtes cuspidales du cône dont nous parlons dépend donc d'une manière symétrique de $[F_1]$ et de $[F_2]$.

On voit ainsi, en se rappelant la symétrie des constructions des deux surfaces auxiliaires $[F_3]$ et $[F_4]$, que \emph{les cônes résidus circonscrits à $[F_3]$ et à $[F_4]$ dont les sommets se trouvent, respectivement, sur les droites multiples $(M_3)$ et $(M_4)$ de ces surfaces, ont les mêmes caractéristiques Plückeriennes}.

A l'endroit déjà cité des Comptes Rendus j'ai démontré cette proposition d'une manière plus immédiate, en prenant, dans la construction de $[F_4]$, pour droite $(M_4)$ une droite qui joint $O_3$ à un point de $(M_3)$ et pour point $O_4$ (point multiple de $[F_4]$) un autre point de $(M_3)$. Alors on voit sans difficulté que le point de rencontre $P$ de $(M_3)$ et de $(M_4)$ est le sommet d'un cône circonscrit à la fois aux deux surfaces auxiliaires $[F_3]$ et $[F_4]$. La proposition démontrée de cette manière permettrait de prouver la formule (II) au moyen d'une seule des deux formules (6) et (7). La formule ($6_c$), par exemple, nous montrerait alors que la classe du cône circonscrit à la fois à toutes les deux surfaces $[F_3]$ et $[F_4]$, est égale à
\[
n'_3 - s = 2\,a_1 + 2\,c_1 - 3\,n_1 + n'_2 + r_2 + \Sigma(\mu_2 + \eta_2 + 2\,\xi_2) + 4\,s
\]
et à
\[
n'_4 - s = 2\,a_2 + 2\,c_2 - 3\,n_2 + n'_1 + r_1 + \Sigma(\mu_1 + \eta_1 + 2\,\xi_1) + 4\,s.
\]

\textbf{26.} \emph{Nombre $c'_3$ des plans tangents stationnaires d'un cône circonscrit à la surface $[F_3]$.} Les plans tangents stationnaires d'un cône circonscrit à une surface sont les plans tangents stationnaires de la surface qui passent par son sommet. Cette série de plans a pour enveloppe une surface développable dont il faut compter les plans tangents

\origpage{39}
passant par un point donné. Prenons le point $O_3$. Les plans cherchés se trouvent alors au nombre des plans tangents à $[F_3]$ qui sortent de $O_3$. Ceux-ci sont (voir le \no 23.): $1^0$ les plans tangents au cône tangent $[O_3\,(S_1)]$ de $[F_3]$ au point $O_3$, $2^0$ les plans qui passent par les droites de $[F_3]$ qui sortent de $O_3$, et $3^0$ les plans tangents au cône $[O_3\,(A_1)]$ circonscrit à $[F_1]$ et au cône $[O_3\,(C_1)]$ projetant son arête cuspidale.

On voit sans difficulté que les plans tangents du cône tangent à un point $s$-tuple d'une surface qui ont un contact stationnaire avec cette surface, seront $1^0$ les plans tangents stationnaires du cône, et $2^0$ ceux dont les génératrices de contact rencontrent la surface à $(s + 2)$ points coïncidant avec le point multiple, à moins que ces génératrices ne soient \emph{seulement} tangentes à des branches de la courbe double et de la courbe cuspidale. Comme toutes les droites le long desquelles ces plans sont tangentes à l'enveloppe des plans tangents stationnaires de la surface, passent par le point multiple, les plans comptent, au moins, pour \emph{deux} au nombre cherché des plans tangents stationnaires qui sortent du point multiple. En regardant un exemple${}^{*)}$ qui n'offre aucune circonstance particulière qui pourrait altérer la règle générale, on voit que les plans de la première classe comptent pour \emph{deux}, et ceux de la seconde classe pour \emph{trois} plans tangents stationnaires de la surface. On trouve ainsi les deux termes suivants${}^{**)}$ du nombre cherché $c'_3$
\[
2\cdot n_{s, 1} + 3\cdot [s(s + 1) - 2\,h_{s, 1}].
\]

Comme les droites de $[F_3]$ qui sortent du point $O_3$ sont au nombre des $[s(s + 1) - 2\,h_{s, 1}]$ droites qui rencontrent $[F_3]$ en $s + 2$ points coïncidant avec $O_3$, les plans qui ont ces droites pour génératrices de contact avec le cône tangent à $O_3$, se trouvent au nombre des plans que nous venons de nommer. Les autres plans par ces droites qui ont des contacts stationnaires avec $[F_3]$ sont, suivant l'article précédent${}^{***)}$, leurs $\bar{\beta}' (= 1\bar{\beta}' + 2\cdot 2\bar{\jmath} + 3\cdot 3\bar{\jmath})$ plans stationnaires pris \emph{trois fois}, et \emph{encore deux fois} les $2\bar{\jmath}$ plans tangents à leurs points-pinces doubles et \emph{quatre fois} les $3\bar{\jmath}$ plans tangents à leurs points-pinces triples, le nombre $\bar{i}$ étant égal à zéro pour toutes les droites.

\textbf{*)} On peut regarder une surface du quatrième ordre douée d'un point triple. --- Les applications de la formule que nous cherchons (voir dans ce qui suit les applications 2 et 3) le montreraient aussi, si nous avions attribué des coefficients indéterminés aux termes dus à ces deux classes de plans.

\textbf{**)} Pour le nombre des plans de la seconde classe voir les \nos 17. et 12.

\textbf{***)} Voir dans le tableau du \no 10. le nombre des plans tangents stationnaires du cône circonscrit résidu.

\origpage{40}
Pour les droites $(D_{1, 3})$ qui joignent $O_3$ aux points fondamentaux $D_1$ de $[F_1]$ on a (voir le \no 16. du présent article)
\[
1\bar{\beta}' = \nu_1 + z_1, \quad 2\bar{\jmath} = \eta_1, \quad 3\bar{\jmath} = \xi_1, \quad \bar{\beta}' = x_1 + z_1{}^{*)},
\]
et pour les autres $n_2 + s$ droites ces nombres caractéristiques sont égaux à zéro. On trouve ainsi le terme suivant du nombre cherché $c'_3$
\[
\Sigma(3\,x_1 + 3\,z_1 + 2\,\eta_1 + 4\,\xi_1).
\]

Un plan tangent à un des cônes $[O_3\,(A_1)]$ et $[O_3\,(C_1)]$ sera plan tangent stationnaire de $[F_3]$ si sa courbe d'intersection avec $[F_3]$ a un point cuspidal. Comme les sections que les plans par $O_3$ font à la surface $[F_3]$ correspondent à celles qu'ils font à la surface $[F_1]$ (voir le \no 23.), les plans cherchés seront $1^0$ les $c'_1$ plans tangents stationnaires de la surface $[F_1]$ (voir le \no 8.), $2^0$ les plans osculateurs à la courbe cuspidale $(C_1)$ de $[F_1]$, dont nous désignerons le nombre par $n_{c, 1}$ (voir le \no 5.), et $3^0$ les $\sigma_1$ plans qui sont tangents à la surface $[F_1]$ à des points de la courbe $(C_1)$ (voir le \no 5.). La tangente de rebroussement de la section que fait un de ces derniers plans, sera la droite qui joint son point cuspidal à $O_3$; car cette droite, qui rencontre $[F_1]$ en trois points consécutifs, doit aussi rencontrer $[F_3]$ en trois points consécutifs${}^{**)}$. Or la tangente de rebroussement de la section faite par un plan tangent stationnaire est génératrice de l'enveloppe des plans tangents stationnaires. Donc chacun des $\sigma_1$ plans dont nous parlons compte pour \emph{deux} plans tangents stationnaires de $[F_3]$ sortant de $O_3$. On trouve ainsi les termes suivants de $c'_3$
\[
c'_1 + n_{c, 1} + 2\,\sigma_1.
\]

L'expression totale de $c'_3$ devient
\[
\begin{aligned}
c'_3 &= 2\cdot n_{s, 1} + 3\cdot [s(s + 1) - 2\,h_{s, 1}] \\
&\quad + \Sigma(3\,x_1 + 3\,z_1 + 2\,\eta_1 + 4\,\xi_1) + c'_1 + n_{c, 1} + 2\,\sigma_1.
\end{aligned} \tag{$8_a$}
\]
En remplaçant $n_{c, 1}$ par l'expression qu'en donnent les relations Plückeriennes
\[
n_{c, 1} = \beta_1 + 3\,r_1 - 3\,c_1,
\]
et $n_{s, 1}$ et $h_{s, 1}$ par les expressions (3), on trouve
\[
\begin{aligned}
c'_3 &= c'_1 + \beta_1 + 3\,r_1 - 3\,c_1 + 2\,\sigma_1 \\
&\quad + \Sigma(3\,x_1 + 3\,z_1 + 2\,\eta_1 + 4\,\xi_1) + 18(p_2 - 1) + 18\,s.
\end{aligned} \tag{$8_b$}
\]

\textbf{27.} On trouve, au moyen du précédent article et du \no 18. du présent article, qu'un cône circonscrit résidu, mené à la surface $[F_3]$

\textbf{*)} Dans la formule (VI) de l'introduction du précédent article nous avons posé $v + 2\,\eta + 3\,\xi = x$.

\textbf{**)} On voit aussi sans difficulté que cette droite qui joint $O_3$ au point cuspidal de la section, que nous appelons $P_{1, 3}$, est tangente à la direction $(P_{1, 3} P'_{1, 3})$ qui est tangente aux courbes par $P_{1, 3}$ et sur $[F_3]$ dont les courbes correspondantes sur $[F_1]$ ont des points cuspidaux (voir les \nos 4. et 5.).

\origpage{41}
d'un point de sa droite singulière $(M_3)$ est doué de $c'_3 - 3\,r_{s, 2}$ plans tangents stationnaires. Comme $c'_3$ a la valeur que nous venons d'indiquer, et $3\,r_{s, 2}$, celle qui est exprimée dans la première formule $(3_b)$, le cône circonscrit résidu aura
\[
\begin{aligned}
& c'_1 + \beta_1 + 3\,r_1 - 3\,c_1 + 2\,\sigma_1 - 6\,(p_1 - 1) \\
& + \Sigma(3\,x_1 + 3\,z_1 + 2\,\eta_1 + 4\,\xi_1) + 18\,(p_2 - 1) + 12\,s
\end{aligned}
\]
plans tangents stationnaires.

Ce nombre doit, selon le \no 25., être égal à celui des plans tangents stationnaires d'un cône circonscrit résidu mené d'un point de la droite $(M_4)$ à l'autre surface auxiliaire $[F_4]$, ou à
\[
\begin{aligned}
& c'_2 + \beta_2 + 3\,r_2 - 3\,c_2 + 2\,\sigma_2 - 6\,(p_2 - 1) \\
& + \Sigma(3\,x_2 + 3\,z_2 + 2\,\eta_2 + 4\,\xi_2) + 18\,(p_1 - 1) + 12\,s.
\end{aligned}
\]
On trouve ainsi l'équation suivante, où nous avons substitué à $2\,(p - 1)$ son expression $a + c - 2\,n$ [voir les formules (1)],
\[
\tag{III}
\left\{
\begin{aligned}
& c'_1 - 12\,a_1 + 24\,n_1 + \beta_1 + 3\,r_1 - 15\,c_1 + 2\,\sigma_1 + \Sigma(3\,x_1 + 3\,z_1 + 2\,\eta_1 + 4\,\xi_1) = \\
& c'_2 - 12\,a_2 + 24\,n_2 + \beta_2 + 3\,r_2 - 15\,c_2 + 2\,\sigma_2 + \Sigma(3\,x_2 + 3\,z_2 + 2\,\eta_2 + 4\,\xi_2).
\end{aligned}
\right.
\]
Les deux membres de cette équation feront, si l'on en soustrait $24$, respectivement pour les deux surfaces $[F_1]$ et $[F_2]$, le $24$-tuple de l'expression que M.~Cayley${}^{*)}$ a donné du genre d'une surface. L'équation (III) exprime donc le théorème de M.~Clebsch${}^{**)}$: \emph{que deux surfaces dont les points se correspondent un-à-un sont du même genre}.

\section*{V. Extensions des résultats exprimés par les équations (I), (II) et (III).}

\textbf{28.} Nous avons supposé dans la déduction des équations (I), (II) et (III) (voir les \nos 12, 24 et 28) qu'aucune des surfaces $[F_1]$ et $[F_2]$ ne soit douée de \emph{points-clos}, c'est à dire de points non-coniques où la courbe cuspidale est rencontrée par les courbes de contact de tous les cônes circonscrits à la surface. S'il y en a sur nos surfaces et \emph{si tous ces points sont fondamentaux} on verra, en y appliquant les mêmes procédés qu'aux autres points fondamentaux, qu'on peut, dans les formules (II) et (III), regarder ces points comme des points coniques doués des caractéristiques suivantes: $\mu = 2$, $v = z = 1$, $y = \eta = \xi = 0$, ($x = 1$). Il n'en

\textbf{*)} Philosophical Transactions 1869, p.~227. Il faut qu'on se rappelle que nous avons supposé que $\chi = \vartheta = 0$ (voir le \no 7. du présent article et l'introduction du précédent article. $\vartheta$ est le nombre des ,,off-points``). Quant à $\chi$, nous y aurons égard dans le \no suivant, et son coefficient sera d'accord avec l'expression de M.~Cayley. Dans ce qui suit nous aurons aussi égard à des points coniques non-fondamentaux.

\textbf{**)} Comptes Rendus de l'Académie des Sciences 1868, Tome LXVII, p.~1238.

\origpage{42}
sera pas de même des formules (I); mais en appliquant les mêmes procédés qui ont donné ces formules à des surfaces douées de points-clos fondamentaux, et en observant qu'une courbe passant par un point-clos y est tangente à toute «première polaire par rapport à la surface», on trouve sans difficulté les modifications que ces formules subiront. $\chi_1$ et $\chi_2$ sont les nombres des points-clos, et nous désignons par $\chi_{1, 2}$ et $\chi_{2, 1}$ les ordres des courbes composées des courbes correspondant à ces nouveaux points fondamentaux. On aura donc, en séparant dans les équations nommées ces points des autres points fondamentaux (et en remplaçant en même temps dans les équations (I) $2\,(p - 1)$ par son expression $a + c - 2\,n$) les équations suivantes:
\[
3\,n_2 - a_2 - c_2 + s\,(n_1 - 4) - \chi_{1, 2} - \Sigma[\mu_{1, 2}\,(\mu_1 - 2)] - b_{1, 2} - 2\,c_{1, 2} = 0, \tag*{[$\mathrm{I}_a$]}
\]
\[
3\,n_1 - a_1 - c_1 + s\,(n_2 - 4) - \chi_{2, 1} - \Sigma[\mu_{2, 1}\,(\mu_2 - 2)] - b_{2, 1} - 2\,c_{2, 1} = 0, \tag*{[$\mathrm{I}_b$]}
\]
\[
\begin{aligned}
n'_1 - 2\,a_1 + 3\,n_1 + r_1 - 2\,c_1 + 2\,\chi_1 + \Sigma(\mu_1 + \eta_1 + \xi_1) \\
= n'_2 - 2\,a_2 + 3\,n_2 + r_2 - 2\,c_2 + 2\,\chi_2 + \Sigma(\mu_2 + \eta_2 + \xi_2),
\end{aligned} \tag*{[II]}
\]
\[
\left\{
\begin{aligned}
& c'_1 - 12\,a_1 + 24\,n_1 + \beta_1 + 3\,r_1 - 15\,c_1 + 2\,\sigma_1 + 6\,\chi_1 \\
&\quad + \Sigma(3\,x_1 + 3\,z_1 + 2\,\eta_1 + 4\,\xi_1) \\
&= c'_2 - 12\,a_2 + 24\,n_2 + \beta_2 + 3\,r_2 - 15\,c_2 + 2\,\sigma_2 + 6\,\chi_2 \\
&\quad + \Sigma(3\,x_2 + 3\,z_2 + 2\,\eta_2 + 4\,\xi_2),
\end{aligned}
\right. \tag*{[III]}
\]
où les sommes $\Sigma$ \emph{sont étendues aux points fondamentaux qui se trouvent aux points coniques et à des points simples}. $b_{1, 2}$, $c_{1, 2}$ et $\mu_{1, 2}$ désignent les ordres des courbes de $[F_2]$ qui correspondent à la courbe double, à la courbe cuspidale et à un point fondamental de $[F_1]$; $b_{2, 1}$, $c_{2, 1}$ et $\mu_{2, 1}$ ont les significations analogues.

\textbf{29.} Nous avons supposé dans le \no précédent que tous les points coniques et tous les points-clos soient fondamentaux; mais les formules trouvées resteront justes, s'il y a sur les deux surfaces:

$1^0$ des points coniques qui se correspondent et qui sont du même ordre $\mu$, du même genre $\pi$ (car $2\,(\pi - 1) = v + z + \xi - 2\,\mu$), et doués des mêmes nombres $\eta$ et $\xi$ de génératrices doubles et stationnaires \emph{non-tangentes} à la courbe double et à la courbe cuspidale,

$2^0$ des points-clos qui se correspondent,

et $3^0$ des points coniques doués des caractéristiques $\mu = 2$, $v = 2$, $y = \eta = \text{etc.} = 0$, et des points-clos qui y correspondent,

\emph{et si la somme $\Sigma$ est étendue à tous les points coniques des surfaces et aux points fondamentaux simples}${}^{*)}$. Comme les points fondamentaux simples n'entrent pas dans l'équation [III], celle-ci exprime

\textbf{*)} Les équations de l'introduction du précédent article ne seront pas altérées si l'on étend la somme $\Sigma$ aussi à des points simples. --- Pour un point conique non-fondamental est $\mu_{1, 2} = 0$ (ou $\mu_{2, 1} = 0$).

\origpage{43}
\emph{une condition nécessaire, mais non pas suffisante}, à laquelle les nombres caractéristiques des deux surfaces doivent satisfaire.

On peut prouver que si deux surfaces dont les points se correspondent un-à-un ont des points coniques qui se correspondent sans satisfaire aux conditions nommées, l'une ou l'autre de ces surfaces aura, au point conique qui s'y trouve, un point fondamental doué d'une direction fondamentale (voir le \no 7.). Les formules [I], [II] et [III] seront donc applicables à deux surfaces, douées des singularités nommées dans l'introduction du précédent article, dont les points se correspondent un-à-un, \emph{si elles n'ont pas}

$1^0$ des courbes cuspidales correspondant l'une à l'autre (voir le \no 3.),

$2^0$ des points fondamentaux doués de directions fondamentales (voir le \no 7.),

$3^0$ des points fondamentaux qui se trouvent à des points-pinces ou à des points de la courbe cuspidale qui ne sont ni points coniques ni points-clos (voir le \no 7.).

\section*{VI. Applications.}

\textbf{1.} \emph{Les surfaces $[F_1]$ et $[F_2]$ sont des plans}${}^{*)}$. Dans ce cas les formules [I] nous donnent
\[
\Sigma(\mu_{1, 2}) = \Sigma(\mu_{2, 1}) = 3\,(s - 1).
\]
$\Sigma(\mu_{1, 2})$ et $\Sigma(\mu_{2, 1})$ sont les ordres des courbes qui, dans chacun des plans, sont composées de toutes les courbes correspondant à des points fondamentaux.

L'équation [II] nous dit que les deux figures planes ont le même nombre de points fondamentaux.

L'équation [III] devient identique.

\textbf{2.} $[F_2]$ \emph{est un plan}, $[F_1]$ \emph{n'a aucune courbe cuspidale}. L'équation [III] nous donne la condition \emph{nécessaire} (mais non pas suffisante):
\[
c'_1 - 12\,a_1 + 24\,n_1 + \Sigma(3\,x_1 + 2\,\eta_1 + 4\,\xi_1) - 24 = 0, \tag{9}
\]
qui exprime que le genre de la surface $[F_1]$ doit être égal à zéro.

L'équation $[\mathrm{I}_a]$ nous donne l'expression suivante de l'ordre de la courbe qui correspond à la courbe double de $[F_1]$
\[
b_{1, 2} = s\,(n_1 - 4) + 3 - \Sigma[\mu_{1, 2}\,(\mu_1 - 2)]. \tag{10}
\]
L'équation $[\mathrm{I}_b]$ nous donne
\[
\Sigma(\mu_{2, 1}) = 3\,s - 3\,n_1 + a_1. \tag{11}
\]
L'équation [II] nous donne l'expression suivante du nombre $m_2$ des points fondamentaux du plan $[F_2]$:

\textbf{*)} Les résultats indiqués ici sont trouvés par M.~Cremona (\emph{Sulle trasformazioni geometriche delle figure piane}, Nota $\mathrm{II}^a$; tomo~V (serie $2^a$) delle \emph{Memorie dell' Academia di Bologna} 1865 et \emph{Giornale di matematiche} tomo~III.

\origpage{44}
\[
m_2 = n'_1 - 2\,a_1 + 3\,n_1 + \Sigma(\mu_1 + \eta_1 + 2\,\xi_1) - 3. \tag{12}
\]
Dans le cas où la surface $[F_1]$ n'a aucun point fondamental les formules (10) et (11) sont démontrées par M.~Clebsch${}^{*)}$, et la formule (12) est d'accord avec les résultats trouvés par ce savant pour des surfaces particulières.

On peut appliquer les formules (9)---(12) au cas d'une surface $[F_1]$ douée d'un point conique $D$, dont nous désignerons les nombres caractéristiques comme ordinairement par $\mu$, $v$ etc.${}^{**)}$, et qui est elle-même de l'ordre $\mu + 1$.

La courbe double de cette surface est composée de $y$ droites doubles, la courbe cuspidale de $z$ droites cuspidales, et il y a encore $\mu\,(\mu + 1) - 4\,y - 6\,z$ droites simples qui se trouvent sur la surface. Toutes ces droites passent par $D$. Il y a sur chacune des $z$ droites cuspidales encore un point conique
\[
(\mu = 3, \quad v = 3, \quad z = 1, \quad y = \eta = \xi = 0).
\]
On peut trouver la classe $n'$ de cette surface et le nombre $c'$ de ses plans tangents stationnaires qui sortent d'un point quelconque, par le même procédé qui nous a donné les mêmes nombres pour la surface $[F_3]$, en cherchant les plans tangents qui passent par le point conique $D$. On trouve les coefficients difficiles à déterminer directement qu'on rencontre dans cette recherche, et les autres nombres caractéristiques de la surface, en faisant usage des formules données dans l'introduction du précédent article. Puis on peut vérifier les résultats trouvés en appliquant les formules (9)---(12) à la surface et à un plan qui y correspond. Les valeurs qu'on trouve sont les suivantes:
\[
\begin{aligned}
n &= \mu + 1 \\
a &= 2\,\mu + x \qquad (x = v + 2\,\eta + 3\,\xi) \\
n' &= 2\,\mu + 3\,x - y - 3\,z - \eta - 2\,\xi \\
c' &= 9\,x - 6\,z - 2\,\eta - 4\,\xi \\
\varkappa &= 3\,x + 3\,y + 3\,z + \eta + 2\,\xi \\
2\,\delta &= x\,(4\,\mu + x - 9) + 6\,z - 2\,\eta - 4\,\xi \\
b &= y, \quad \varrho = y, \quad j = 2\,y, \quad f = 0 \\
c &= z, \quad r = 0, \quad \sigma = 0, \quad \chi = z, \quad \beta = 0, \quad i = 0 \\
&\phantom{{}=} \qquad\qquad\qquad \text{etc.}
\end{aligned}
\]
Nous projetterons cette surface du centre $D$ sur un plan. Alors les points de la surface et du plan se correspondent un-à-un. Afin de

\textbf{*)} \emph{Intorno alla représentazione di superficie algebriche sopra un piano}, Istituto Lombardo 1868; \emph{Mathematische Annalen} Bd.~I, p.~280 etc.

\textbf{**)} Quoique $z$ doive être égal à zéro dans l'application actuelle, je suppose qu'il ne sera pas dénué d'intérêt de connaître les nombres caractéristiques de la surface dans les cas où $z$ n'est pas zéro.

\origpage{45}
pouvoir faire usage des résultats trouvés dans cet article-ci, nous devons supposer que $z = 0$; car à une droite cuspidale correspondrait un point fondamental doué d'une direction fondamentale. Alors on a $c = 0$ [ce qu'on a aussi supposé dans les formules (9)---(12)]. La surface donnée n'a qu'un seul point fondamental, qui est placé au point conique d'ordre $\mu$. La courbe qui y correspond est de l'ordre $\mu_{1, 2} = \mu$. Les points fondamentaux du plan sont ceux où il est rencontré par les droites de la surface, et les courbes correspondantes sont ces droites, dont les $y$ sont doubles ($\mu_{2, 1} = 2$) pendant que les autres $\mu\,(\mu + 1) - 4\,y$ ($= 2\,\mu + x - 2\,y$) sont simples ($\mu_{2, 1} = 1$). On a donc
\[
\begin{aligned}
m_2 &= 2\,\mu + x - y, \\
\Sigma(\mu_{2, 1}) &= 2\,\mu + x.
\end{aligned}
\]
Comme les droites doubles correspondent à des points, on a $b_{1, 2} = 0$. Ces valeurs sont les mêmes qu'on trouverait au moyen des équations (10), (11) et (12), $s$ étant égal à $\mu + 1$.

Les nombres caractéristiques de la surface satisfont aussi à l'équation (9). Les coefficients de $\eta$ et de $\xi$ dans cette équation, qui est un cas particulier de [III], sont déterminés au moyen de cet exemple-ci, et c'est ainsi que nous avons trouvé les coefficients de $2\,\bar{\jmath}$ et de $3\,\bar{\jmath}$ dans l'expression, indiquée dans le \no 10. du précédent article, du nombre des \emph{plans tangents stationnaires du ,,cône circonscrit résidu``}.

\textbf{3.} \emph{Les surfaces $[F_1]$ et $[F_2]$ sont deux surfaces réciproques} ($[F]$ et $[F']$). Dans ce cas on a
\[
\begin{aligned}
n_1 &= n = n'_2, \\
n'_1 &= n' = n_2, \\
a_1 &= a = a' = a_2, \\
&\quad\ \text{etc.}
\end{aligned}
\]
Un point où il y a une infinité de plans tangents sera un point fondamental, et s'il y a, au nombre de ces plans tangents à un même point, une infinité qui passent par une même droite, celle-ci sera une direction fondamentale. Les points fondamentaux se trouveront donc aux points coniques, aux points-pinces et aux points-clos. Ces dernières deux classes de points auront toujours des directions fondamentales, et les points coniques en auront, si les cônes tangents sont doués de génératrices doubles et cuspidales \emph{non-tangentes} à la courbe double et à la courbe cuspidale. Comme nous n'avons pas eu égard, dans la déduction de nos formules, à des directions fondamentales, nous devons supposer qu'il n'y en ait sur aucune des deux surfaces, ou que
\[
j = j' = \chi = \chi' = 0,
\]

\origpage{46}
et que, pour tout point conique et pour tout plan tangent le long d'une courbe de $[F]$,
\[
\eta = \xi = 0 \quad\text{et}\quad \eta' = \xi' = 0,
\]
d'où
\[
x = v \quad\text{et}\quad x' = v'.
\]
La courbe correspondant à un point conique est celle le long de laquelle le plan correspondant est tangent à la surface réciproque, de façon que pour tous ces points et plans
\[
\mu_{1, 2} = v = x \quad\text{et}\quad \mu_{2, 1} = v' = x'.
\]
On aura encore
\[
s = a,
\]
\[
b_{1, 2} = \varrho, \quad b_{2, 1} = \varrho',
\]
\[
c_{1, 2} = \sigma, \quad c_{2, 1} = \sigma'.
\]
L'équation $[\mathrm{I}_a]$ nous donnera alors, si l'on y remplace $3\,n' - c'$ par la quantité $3\,a - \varkappa$, qui y est égale (équation réciproque à la $3^{\text{me}}$ équation (II) de l'introduction du précédent article):
\[
a\,(n - 2) = \varkappa + \varrho + 2\,\sigma + \Sigma[x\,(\mu - 2)],
\]
qui est la $1^{\text{re}}$ équation (VIII) du précédent article.

L'équation $[\mathrm{I}_b]$ donne l'équation réciproque.

Nos deux autres équations, [II] et [III], nous donnent (comme $a = a'$) les résultats suivants
\[
2\,n + r - 2\,c + \Sigma(\mu) = 2\,n' + r' - 2\,c' + \Sigma'(\mu'), \tag{13}
\]
\[
\begin{aligned}
24\,n + \beta + 3\,r - 16\,c + 2\,\sigma + \Sigma(3\,x + 3\,z) \\
= 24\,n' + \beta' + 3\,r' - 16\,c' + 2\,\sigma' + \Sigma'(3\,x' + 3\,z').
\end{aligned} \tag{14}
\]
On ne peut déduire ces résultats-ci des équations qui sont indiquées dans l'introduction du précédent article et des équations réciproques, mais bien l'équation suivante qu'on trouve en soustrayant les deux membres de (14) de ceux de l'équation (13) après avoir multiplié ces derniers par $6$:
\[
\left\{
\begin{aligned}
& - 12\,n + 3\,r + 4\,c - \beta - 2\,\sigma + \Sigma(6\,\mu - 3\,x - 3\,z) = \\
& - 12\,n' + 3\,r' + 4\,c' - \beta' - 2\,\sigma' + \Sigma'(6\,\mu' - 3\,x' - 3\,z').
\end{aligned}
\right. \tag{15}
\]
L'une des deux équations (13) et (14) peut remplacer l'une des deux équations que nous avons negligées dans l'introduction du précédent article.

Au lieu de déduire l'équation (15) des équations du précédent article, au nombre desquelles se trouve l'équation $i' = i$, on peut faire usage de l'équation (15), trouvée ici par une autre voie, pour démontrer${}^{*)}$ que $i' = i$.

\textbf{*)} Voir le \no 11. du précédent article.

\origpage{47}
En effet, on trouve au moyen de la première équation (II)${}^{*)}$ et de l'équation (IV)
\[
c \cdot n\,(n - 1) = a\,c + 2\,b\,c + 3\,c + 3\,r + 6\,h + 9\,\beta.
\]
En en soustrayant la dernière équation (IX) et la dernière équation (VIII) après avoir multiplié celle-ci par quatre, on trouve (comme $\chi = \eta = \xi = 0$)
\[
c + 3\,r - \beta - 5\,\sigma + 2\,i - \Sigma(z) + \Sigma'(2\,v') = 0. \tag{16}
\]
On déduit, de la même manière, de la première équation (II), de l'équation réciproque à la deuxième équation (II) et des premières équations (VIII) et (IX):
\[
- a + n' - \varkappa + \sigma + \Sigma(x) = 0,
\]
ou bien, comme
\[
\varkappa = c' + 3\,a - 3\,n' \quad\text{[équation réciproque à la } 3^{\text{me}} \text{ (II)]},
\]
\[
4\,n' - 4\,a - c' + \sigma + \Sigma(x) = 0. \tag{17}
\]
On voit, au moyen des équations (16) et (17) et des équations réciproques (et de l'équation $a' = a$), que les quantités
\[
c + 3\,r - \beta - 5\,\sigma + 2\,i - \Sigma(z + 2\,v)
\]
et
\[
- 4\,n + c + \sigma + \Sigma(x)
\]
sont égales aux quantités réciproques qu'on forme en substituant les lettres accentuées aux lettres non-accentuées. La quantité
\[
- 12\,n + 4\,c + 3\,r - \beta - 2\,\sigma + 2\,i + \Sigma(3\,x - z - 2\,v),
\]
qui en est composée, satisfait donc à la même condition.

La même chose a lieu, selon l'équation (15), par rapport à la quantité
\[
- 12\,n + 4\,c + 3\,r - \beta - 2\,\sigma + \Sigma(6\,\mu - 3\,x - 3\,z),
\]
et, par conséquent, aussi par rapport à
\[
2\,i + \Sigma(6\,x - 6\,\mu + 2\,z - 2\,v).
\]
Or on a pour chaque point conique de $[F]$ et pour chaque plan tangent à $[F]$ le long d'une courbe [$2^{\text{me}}$ équation (VII)]
\[
\begin{aligned}
v &= z + 3\,x - 3\,\mu, \\
v' &= z' + 3\,x' - 3\,\mu';
\end{aligned}
\]
(car, dans le cas actuel, où $\eta = 0$ et où l'on doit regarder séparément une nappe passant par un point conique formé par d'autres nappes, le cône tangent à un point conique ne peut être composé de parties au nombre desquelles il y a des plans). Donc
\[
i = i'.
\]

\textbf{*)} Les formules de l'introduction du précédent article sont, dans cette démonstration, les seules qui sont marquées de chiffres romains. ---

Nous suivons, dans la déduction des équations (16) et (17), des calculs faits par M.~Cayley. Phil. Trans. 1869, pp.~203 et 227.

\origpage{48}
\section*{VII. Addition sur la représentation d'une surface sur un plan multiple.}

On peut appliquer les mêmes procédés dont nous avons fait usage dans ce mémoire, à la recherche des propriétés de deux surfaces dont les points se correspondent d'une manière \emph{multiple}. Nous nous contenterons ici de nous occuper de la représentation d'une surface sur un plan multiple, disons $N$-tuple. Dans ce cas, à chaque point de la surface correspond un seul point d'un plan, mais à un point du plan correspondent $N$ points de la surface. Pour les points d'une certaine courbe du plan, \emph{la courbe de liaison} (Uebergangscurve), deux de ces $N$ points coïncident. On voit que cette courbe jouera le même rôle que, dans ce qui précède, la courbe de contact d'un cône circonscrit ou la courbe cuspidale, et nous désignerons son ordre par $A$, sa classe par $N'$ et le nombre de ses tangentes stationnaires par $C'$. Le plan multiple peut contenir des points fondamentaux: l'ordre $\mu$ d'un de ces points fondamentaux indique le nombre de nappes du plan multiple qui contribuent à le former [les points des ($N - \mu$) autres nappes qui coïncident avec le point fondamental ont pour correspondants ($N - \mu$) points isolés de la surface], et la classe $v$ indique le nombre des branches de la courbe de liaison qui passent par le point.

Si nous supposons que c'est la surface $[F_2]$ qui se réduit à un plan multiple, nous aurons, en faisant l'usage ordinaire des indices 1 et 2 et de la notation $s$, les équations
\[
3\,N - A + s\,(n_1 - 4) - \chi_{1, 2} - \Sigma[\mu_{1, 2}\,(\mu_1 - 2)] - b_{1, 2} - 2\,c_{1, 2} = 0 \tag*{[1]}
\]
\[
\begin{aligned}
n'_1 - 2\,a_1 + 3\,n_1 + r_1 - 2\,c_1 + 2\,\chi_1 + \Sigma[\mu_1 + \eta_1 + 2\,\xi_1] \\
= N' - 2\,A + 3\,N + \Sigma(\mu_2)
\end{aligned} \tag*{[2]}
\]
\[
\left\{
\begin{aligned}
& c'_1 - 12\,a_1 + 24\,n_1 + \beta_1 + 3\,r_1 - 15\,c_1 + 2\,\sigma_1 + 6\,\chi_1 \\
&\quad + \Sigma(3\,x_1 + 3\,z_1 + 2\,\eta_1 + 4\,\xi_1) = C' - 12\,A + 24\,N + \Sigma(3\,v_2),
\end{aligned}
\right. \tag*{[3]}
\]
qu'on prouve comme les équations analogues $[\mathrm{I}_a]$, [II] et [III]. Dans la démonstration de [1] on fait usage de l'équation suivante, que j'ai prouvée dans les ,,Mathematische Annalen`` 3.~Bd. p.~324,
\[
2\,(P - 1) = A - 2\,N, \tag*{[4]}
\]
où $P$ est le genre de la courbe de $[F_1]$ qui correspond à une section plane du plan multiple. (On pourrait appeler $P$ le genre de cette section qui est une droite multiple. L'analogie de l'équation [4] avec l'équation (1) dans le \no 1 est alors évidente.)

En appliquant au cas actuel les mêmes procédés qui donnent l'équation $[\mathrm{I}_b]$, on voit qu'il faut remplacer, dans cette équation,
\[
(s\,(n_2 - 1) - \Sigma[\mu_{2, 1}\,(\mu_2 - 1)] - b_{2, 1} - 3\,c_{2, 1}) + c_{2, 1}
\]

\origpage{49}
par l'ordre $A_{2, 1}$ de la courbe qui sur $[F_1]$ correspond à la courbe de liaison. On trouve ainsi
\[
A_{2, 1} = a_1 + c_1 - 3\,n_1 + 3\,s - \Sigma(\mu_{2, 1}). \tag*{[5]}
\]
Pour compléter ces résultats nous placerons ici l'équation générale dont l'équation [4] est un cas particulier:
\[
A = 2\,(P_{2, 1} - 1) - 2\,N\,(P_2 - 1),
\]
où $P_2$ est le genre d'une courbe dans le plan d'image pénétrant toutes ses $N$ nappes, et $P_{2, 1}$ celui de la courbe correspondante. Cette équation est prouvée à l'endroit cité.

M.~Clebsch a appliqué la représentation sur un plan double à la recherche des propriétés \emph{d'une surface du cinquième ordre douée d'une courbe double du quatrième ordre de la première espèce}${}^{*)}$. Cette surface a les nombres caractéristiques:
\[
\begin{gathered}
n = 5, \quad a = 12, \quad c = r = \sigma = \chi = \beta = 0 \\
n' = 20, \quad c' = 48.
\end{gathered}
\]
Dans la représentation de cette surface sur un plan double dont M.~Clebsch fait usage on a
\[
\begin{gathered}
N = 2, \quad s = 5, \quad A = 6, \quad N' = 16, \quad C' = 30, \\
b_{1, 2} = 8, \quad A_{2, 1} = 7, \quad P = 2.
\end{gathered}
\]
Il y a sur la surface trois points fondamentaux simples, auxquels correspondent des droites ($\mu_1 = 1$, $\mu_{1, 2} = 1$). Dans le plan il y a deux points fondamentaux: on a
\[
\begin{aligned}
\text{pour l'un} \quad &\mu_2 = 2, \quad v_2 = 4, \quad \mu_{2, 1} = 3, \\
\text{pour l'autre} \quad &\mu_2 = 2, \quad v_2 = 2, \quad \mu_{2, 1} = 2.
\end{aligned}
\]
Ces quantités satisfont aux équations [1]---[5], qui pourraient aider à les trouver. ---

Nous proposons comme exercice l'application de nos formules à la question suivante:

\emph{Qu'on représente une surface du sixième ordre, douée de deux droites doubles qui ne se rencontrent pas, sur un plan double en projetant tous ses points sur le plan par des droites qui rencontrent les droites doubles; quelles seront alors les propriétés de l'image, et quelles propriétés de la surface peut-on trouver au moyen de cette représentation?}

Quant à la solution de ces questions je me suis borné, moi-même, à la détermination des nombres qui entrent dans mes formules.

\textbf{*)} Abhandlungen der Gesellschaft der Wissenschaften zu Göttingen 15.~Bd.

Copenhague, Février 1871.

\end{document}
